% Archived RhoCalc section from CeTTa.tex, extracted 2026-06-23, pre full rewrite.
% Reference only; not part of the build.

\section{\RhoCalc{}: A Concurrent Reaction Machine}
\label{sec:rhocalc}

So far \CeTTa{} has been a MeTTa evaluator: a term goes in, normal forms come
out.  Under \code{--lang rhocalc} the same runtime becomes something different in
kind---a \emph{concurrent process engine} whose unit of work is not a rewrite but
a \emph{reaction}: a rendezvous between a sender and a receiver on a shared
channel.  This gives MeTTa three things plain evaluation lacks---channels and
private names, genuine concurrency, and, the reason it earns its own section,
\emph{evaluation triggered by communication}.  That last capability, together
with the MeTTa library that exposes it, is what we call \emph{rhometta}.

\subsection{A reaction, in one picture}

The calculus is Meredith and Radestock's reflective higher-order
$\rho$-calculus~\cite{rho2005}.  Its single rewrite rule, \textsc{comm}, is a
handshake: a sender \code{x!(m)} offering a message \code{m} on channel \code{x},
sitting next to a receiver \code{for(y <- x)\{P\}} waiting on that same channel,
react---the message is consumed and the receiver resumes with \code{m}
substituted for its bound name \code{y}:
\[
  x!(m) \;\big|\; \mathsf{for}(y \leftarrow x)\{P\}
  \;\;\xrightarrow{\;\textsc{comm}\;}\;\; P[y := m].
\]

\begin{figure}[h]\centering
\begin{tikzpicture}[>=Stealth,
  box/.style={draw, rounded corners, align=center, inner sep=4pt,
    font=\small\ttfamily}]
  \node[box] (s) {x!(m)};
  \node (par) [right=3mm of s, font=\large] {$\mid$};
  \node[box] (r) [right=3mm of par] {for(y<-x)\{P\}};
  \node[box, fill=black!8] (out) [right=16mm of r] {P[y:=m]};
  \draw[->, thick] (r) -- node[above, font=\scriptsize\sffamily] {\textsc{comm}} (out);
\end{tikzpicture}
\caption{A reaction: a send and a matching receive on channel \code{x} rendezvous;
  the receiver continues with the message substituted in.}
\end{figure}

A little vocabulary makes the rest precise.  \code{0} is the inert process;
\code{P | Q} runs \code{P} and \code{Q} concurrently; \code{new x in P} makes a
fresh private channel.  Two reflection operators let processes become names and
names become processes again: \code{@P} \emph{quotes} a process into a name, and
\code{*x} \emph{drops} (runs) the process a name quotes, with \code{*@P = P}.
Channels are themselves quoted processes---this is what makes the calculus
\emph{higher-order}, and it is the hinge of the next subsection.

\subsection{Eval-at-\textsc{comm}: communication that computes}

Reflection turns an ordinary handshake into a computation.  Let the receiver do
nothing but \emph{drop} whatever it receives:
\[
  \texttt{eval\_drop}(x,p)\;\equiv\;
  \underbrace{x!(@p)}_{\text{send the quoted program } p}\;\big|\;
  \underbrace{\mathsf{for}(y \leftarrow x)\{\,*y\,\}}_{\text{run whatever arrives}}.
\]
The \textsc{comm} delivers the name \code{@p}; the drop \code{*@p} runs \code{p};
the process settles at the \emph{result of running} \code{p}.  A handshake has
become an evaluation step: \emph{communication is computation}.  And because
\code{p} may itself evaluate to \emph{more process}, a reaction can plant the
next reactions---\emph{communication is computation is construction}.  This
eval-at-\textsc{comm} pattern, and its MeTTa library, is rhometta.

\subsection{Engine, profiles, libraries}

rhometta is three layers over the engine, not one monolith (Figure~\ref{fig:rho-layers}).

\begin{figure}[h]\centering
\begin{tikzpicture}[
  lyr/.style={draw, rounded corners, minimum width=86mm, align=center,
    inner sep=4pt, font=\small}]
  \node[lyr] (a) {demos --- \code{benchmarks/rho/rhometta-*} (self-unfolding prover, AGI showcase)};
  \node[lyr, below=1.6mm of a] (b) {\code{lib/rhometta.metta} --- \code{rhometta:run}, \code{:eval}, \code{:transitions}};
  \node[lyr, below=1.6mm of b] (c) {\code{lib/rho.metta} --- strict-core $\rho$ constructor surface};
  \node[lyr, below=1.6mm of c, fill=black!8] (d) {\code{--lang rhocalc} engine \;|\; profiles: \code{strict-core}, \code{cost}};
\end{tikzpicture}
\caption{The rhometta stack.  A \emph{profile} is a named variant of the
  reduction relation; a \emph{library} is MeTTa-level surface that delegates
  reduction to the engine.}
\label{fig:rho-layers}
\end{figure}

\begin{description}
\item[The engine.] \code{--lang rhocalc} reduces $\rho$ terms in either a
  Rholang-like \code{.rho} surface or a MeTTa-style s-expression \code{.mrho}
  surface.  Reduction lives in a dedicated \code{RhoMachine} (run budget,
  fresh-name counter, and a live index of sends and receive-sites rebuilt after
  each reaction) rather than being threaded through the ordinary evaluator.
  Internally allocated fresh names are opaque identities; quoted process names
  are compared as whole names and never bound through, so the fast index covers
  only directly comparable names and quoted names take a structural fallback.
\item[Profiles.] A \emph{profile} names a variant of the reduction relation.
  \code{strict-core} is the pure asynchronous $\rho$ spine; \code{cost} accounts
  the resource each reaction consumes (Section~\ref{sec:rho-cost}).  As
  everywhere in \CeTTa{}, a profile fixes what is \emph{interpreted}, never what
  is legal.
\item[Libraries.] \code{lib/rho.metta} is the base vocabulary for building
  strict-core $\rho$ terms; \code{lib/rhometta.metta} is the rhometta surface
  (\code{rhometta:run} reduces to the quiescent outcome set,
  \code{rhometta:eval} is eval-at-\textsc{comm}, \code{rhometta:transitions}
  gives one-step successors), delegating reduction to the engine.
\end{description}

\subsection{What you observe is what it computes}
\label{sec:rho-lean}

A fair worry: when eval-at-\textsc{comm} fires, is the result the engine reports
really the result of running the payload---or could the machinery quietly change
the answer?  We settle this in Lean.

In plain terms: run \code{eval\_drop(x,p)} to quiescence and collect every
outcome you can observe \emph{up to structural rearrangement} of the process.
That set of observations is exactly the set of results a correct evaluation of
\code{p} would produce---no more, no less---and the statement refers only to what
is observable, never to any internal ``this is the answer'' tag.

Precisely, the theorem \code{rhometta\_bridge\_scObserved} (in
\code{RhometaReduction.lean}) equates the concrete quiescent run of
\code{eval\_drop(x,p)} with \code{RhometaSCObservedOutcomes} of the same
process---the outcomes visible up to structural congruence---with no
result-certificate carried in the answer.  It is axiom-clean
(\code{\{propext, Classical.choice, Quot.sound\}}, zero sorries).  Two ingredients
are load-bearing and worth stating honestly:
\begin{itemize}
\item \textbf{The proof technique.}  Inductions directly on structural
  congruence stall at the symmetry and transitivity steps.  The fix is a
  canonical-form normalizer \code{nfAtom} that is \emph{complete}: two processes
  are structurally congruent exactly when they share a normal form.  Equality of
  normal forms is then an ordinary equation, and the stall is gone.
\item \textbf{One side condition.}  The communicated variable must not occur
  \emph{under a quote} (a hygiene predicate, \code{noBoundUnderQuote}).  This is
  necessary, not bureaucratic: a payload that hides a bound occurrence inside a
  quoted name can defeat the characterization, and a concrete counterexample
  exhibits exactly that failure.
\end{itemize}

\subsection{rhometta as a cognitive substrate}

eval-at-\textsc{comm} is more than a soundness curiosity---it is a way to
program.  Because a reaction's payload can emit new reactions, a reaction network
can \emph{grow itself}.  Two self-contained demos (both run on the
\code{make BUILD=core} binary, under \code{benchmarks/rho/rhometta-*}) show what
that buys.

\paragraph{A self-unfolding prover.}  Seed a \emph{single} cell for a goal.  Each
\textsc{comm} evaluates one backward-chaining step whose result is the cells for
that goal's subgoals, so \code{rhometta:run} keeps reducing and the reaction
network unfolds itself into a proof tree.  Where a goal has several rules the
payload has several results, so the \textsc{comm} \emph{forks}: nondeterministic
inference becomes branching reduction, and \code{rhometta:run} returns one
quiescent residual per complete proof tree (Figure~\ref{fig:cake})---two for the
sample ``can I bake a cake?'' goal, while a single-rule goal stays confluent at
one.

\begin{figure}[h]\centering
\begin{tikzpicture}[>=Stealth, level distance=8mm, sibling distance=15mm,
  every node/.style={font=\scriptsize\ttfamily}, edge from parent/.style={draw,-}]
  \begin{scope}
  \node {cake} child { node {batter} child { node {flour} } child { node {egg} } }
               child { node {oven} };
  \end{scope}
  \begin{scope}[xshift=42mm]
  \node {cake} child { node {batter} child { node {mix} } }
               child { node {oven} };
  \end{scope}
\end{tikzpicture}
\caption{The self-unfolding prover's may-set for ``cake'': two complete proof
  trees, one per rule for \code{batter}.  The trees are not enumerated up front---%
  they are \emph{grown} by reactions whose payloads plant their children.}
\label{fig:cake}
\end{figure}

\paragraph{Cognitive services as parallel reactions.}  A larger showcase runs
three OpenCog/Hyperon primitives concurrently---a PLN two-hop proof search,
overlap-corrected evidence revision across two agents, and one ECAN
attention-spreading round---and chains their channel-gathered results into a
final payload that synthesizes one action plan per proof path.  Plain Rholang can
pass names and spawn processes but cannot evaluate a MeTTa term at \textsc{comm};
a bare MeTTa interpreter has the inference but no channels, concurrency, or
may-set.  rhometta is the fusion, and this showcase is the smallest program that
exhibits all of it at once.

\paragraph{Backward chaining, in the standard idiom.}  The cake prover is the
intuition; a small ladder of demos puts the reaction machine in the exact
representation the Hyperon community recognizes as backward chaining---typed judgments
\code{(: proof theorem)}, rules \code{(: r (-> (: p1 P1) (: p2 P2) Concl))}, and
Curry--Howard proof terms---and they form a deliberate ladder
(Figure~\ref{fig:rho-ladder}).  In \code{rhometta-typed-bc} the recognizable mortal
syllogism (\code{researcher}\,$\wedge$\,\code{curious}\,$\Rightarrow$\,%
\code{mortal}) is proved by reactions: the two premises run as parallel cells, and
the join re-matches the rule so its shared variable is unified---a spurious pairing
(\code{researcher}~$r_1$ with \code{curious}~$e_1$) cannot survive.  In
\code{rhometta-type-inference} type inference \emph{is} the proof search: an
overloaded \code{plus} produces a genuine may-set (an \code{Int} and a \code{Float}
typing of \code{(plus one one)}), while \code{(plus x one)} with \code{x\,:\,Int}
prunes the \code{Float} overload at the unification join.  Proof terms stay pure,
\code{(: (app pf pa pb) (HasType e t))}; a PLN truth value rides \emph{outside} the
judgment as derivation metadata that ranks the overloads.  Each demo is gated by
negative canaries---a bad pairing yields no proof and, because the join is made
total, no stuck reaction---and is checked differentially against a sequential nilbc
chainer on the same knowledge base, which produces the same proof-term skeletons.
Two companion type-inference demos then separate the two deeper motifs:
\code{rhometta-type-inference-nested} sends a deterministic sub-proof back to its
parent join, while \code{rhometta-type-inference-mayset} uses
\code{rho:patterns:!recv} to let two inner typings feed one outer continuation.

\begin{figure}[h]\centering
\begin{tikzpicture}[>=Stealth, node distance=3mm,
  rung/.style={draw, rounded corners, align=left, inner sep=4pt, font=\footnotesize,
    text width=104mm},
  hero/.style={rung, line width=1pt, draw=orange!85!black, fill=orange!8}]
  \node[rung] (a) {\textbf{\code{rhometta-typed-bc}} --- standard nilbc shape.\\
    {\scriptsize\itshape facts/rules $\to$ Curry--Howard proof terms;\quad output: 2 mortal proofs.}};
  \node[rung, below=of a] (b) {\textbf{\code{rhometta-type-inference}} --- proof search as typing.\\
    {\scriptsize\itshape overloads fork, the join prunes bad types;\quad output: Int/Float may-set, then Int-only.}};
  \node[rung, below=of b] (n) {\textbf{\code{rhometta-type-inference-nested}} --- recursive sub-proof flow.\\
    {\scriptsize\itshape a deterministic sub-proof returns to its parent join;\quad output: \code{(f (g x))} : Bool.}};
  \node[rung, below=of n] (m) {\textbf{\code{rhometta-type-inference-mayset}} --- recursive may-set at depth.\\
    {\scriptsize\itshape derived inner proofs fan in via \code{!recv};\quad output: boxed typings, including multi-premise.}};
  \node[hero, below=of m] (c) {\textbf{\code{rhometta-deployment-planner}} --- application-shaped constraint join.\\
    {\scriptsize\itshape 7 premises in parallel, shared variables agree at the join;\quad output: 2 valid EU chat plans.}};
  \node[rung, below=of c] (d) {\textbf{\code{rhometta-agi-showcase}} --- cognitive workflow.\\
    {\scriptsize\itshape PLN $+$ belief revision $+$ ECAN $\to$ plans;\quad output: one plan per proof path.}};
  \draw[->] (a)--(b); \draw[->] (b)--(n); \draw[->] (n)--(m); \draw[->] (m)--(c); \draw[->] (c)--(d);
\end{tikzpicture}
\caption{The rhometta demo ladder, from the smallest recognizable
  backward-chaining shape up to an AGI cognitive workflow.  The first two isolate
  the proof-search motifs; the next two add recursion (deterministic, then a
  may-set at depth via replicated input); the highlighted deployment planner is
  the reader-facing, application-shaped showcase.}
\label{fig:rho-ladder}
\end{figure}

The first two demos are intentionally small: they isolate the two proof-search
motifs---a shared-variable AND-join, and may-set OR-forking with join-time pruning.
The third, \code{rhometta-deployment-planner}, composes those motifs into an
application-shaped query (Figure~\ref{fig:rho-join}).  A deployment plan is not
selected by post-filtering a list; it is \emph{derived} by backward chaining: rho
supplies the concurrent premise cells, and the join enforces that the same request,
model, GPU, provider, region, and latency are used across all seven premises.  A
plan for \code{eu-chat} (an EU, low-latency, GDPR-safe dialogue request) exists only
if all seven judgments agree---request task/privacy/offer, model task/privacy/%
GPU-need, and a matching provider offer---so the engine returns the two valid plans
(\code{cloudA}+\code{llm70bq4} on H100 and \code{cloudB}+\code{llm7b} on L40S) while
the US-region, standard-latency, and HIPAA-only candidates die at the join.

\begin{figure}[t]\centering
\resizebox{\textwidth}{!}{%
\begin{tikzpicture}[>=Stealth, font=\scriptsize,
  req/.style={draw, rounded corners, align=left, inner sep=4pt, fill=black!4},
  cell/.style={draw, rounded corners, fill=blue!8, inner sep=2.5pt, align=left,
    font=\scriptsize, text width=46mm},
  jbox/.style={draw, line width=1pt, draw=orange!85!black, fill=orange!12,
    align=center, inner sep=5pt, rounded corners, text width=26mm},
  good/.style={draw, fill=green!14, align=left, inner sep=4pt, rounded corners, text width=32mm},
  bad/.style={draw, fill=black!7, align=left, inner sep=4pt, rounded corners,
    text=black!55, text width=32mm}]
  \node[req] (req) {\textbf{Request} \code{eu-chat}\\ task $=$ Dialogue\\
    privacy $=$ GDPR\\ region $=$ EU\\ latency $=$ LowLatency};
  \node[cell] (p1) [above right=14mm and 8mm of req] {ReqTask(req, task)};
  \node[cell, below=1mm of p1] (p2) {ReqPrivacy(req, privacy)};
  \node[cell, below=1mm of p2] (p3) {ModelTask(model, task)};
  \node[cell, below=1mm of p3] (p4) {ModelPrivacy(model, privacy)};
  \node[cell, below=1mm of p4] (p5) {ModelNeeds(model, gpu)};
  \node[cell, below=1mm of p5] (p6) {ProviderOffer(provider, gpu, region, latency)};
  \node[cell, below=1mm of p6] (p7) {ReqOffer(req, region, latency)};
  \node[jbox] (join) [right=10mm of p4] {\textbf{JOIN}\\ shared variables agree:
    req, model, provider, gpu, region, latency};
  \node[good] (good) [above right=1mm and 10mm of join] {\textbf{Valid may-set}\\
    cloudA $+$ llm70bq4 (H100)\\ cloudB $+$ llm7b (L40S)};
  \node[bad] (bad) [below right=1mm and 10mm of join] {Pruned\\
    cloudC: wrong region (US)\\ cloudD: wrong latency\\ voicegen: wrong privacy};
  \draw[->, thick] (req) -- (p4.west);
  \foreach \p in {p1,p2,p3,p4,p5,p6,p7} \draw[->, blue!55] (\p.east) -- (join.west);
  \draw[->, green!55!black, thick] (join) -- (good);
  \draw[->, black!45, thick] (join) -- (bad);
\end{tikzpicture}}
\caption{The deployment planner's constraint join.  Each premise is solved in its
  own rho cell (blue); independent premise search \emph{generates} invalid
  cross-products, but they are killed at the join (gold), where the shared variables
  must agree.  The surviving residuals are exactly the valid plans (green); the
  pruned branches (grey) never reach the may-set.}
\label{fig:rho-join}
\end{figure}

The ladder keeps the reader-facing story staged.  The deployment planner prunes a
premise \emph{cross-product} at the join, so it is a semantic demo rather than an
indexed large-catalog planner: a production planner would thread bindings or index
premise lookup earlier.

Two companion demos carry the deeper recursive motifs.
\code{rhometta-type-inference-nested} exercises a \emph{deterministic} recursive
search with a join: the sub-proof for \code{(g x)} flows back along a channel to
the parent reaction for \code{(f (g x))}, which assembles the nested proof term.
\code{rhometta-type-inference-mayset} handles the sharper case---a recursive
may-set \emph{at depth}, where a compound subgoal has several proofs that must all
reach one parent continuation---and here the mechanics matter.  The inner subgoal
\code{(box one)} is type-checked by a real overload \textsc{comm}-fork over the
knowledge base: both \code{(Boxed Int)} and \code{(Boxed Float)} are
\emph{derived}, their proof terms are built by re-matching the rule against the
resolved premise, and the result is cross-checked by set membership against an
independent sequential nilbc chainer rather than against a constant.  One boundary
is stated plainly rather than hidden.  Because the engine forks a multi-result
payload into separate may-set \emph{branches}---one delivered message per branch---a
single send would give a replicated input nothing to gather; so the derived inner
set is collapsed to data, \emph{re-lifted} as parallel sends on the parent channel,
and \code{rho:patterns:!recv} fans them in.  This ``derived-then-delivered'' bridge
means the inner search and the outer join are two reductions joined by a collapse
and re-lift, not one uninterrupted rho reduction.

The fan-in itself is genuine: \code{rho:patterns:!recv} is the lazy
Meredith--Radestock replicated-input encoding, derived from strict-core reflection
(quote/drop), so a primitive persistent receive would be a performance and
residual-cleanliness convenience, not an expressiveness requirement.  Nor is
multiplicity at depth confined to a single overloaded premise: a further check
derives \code{(box (plus one one))}---an eight-way premise cross-product
($2\times2\times2$ over an overloaded operator and two overloaded arguments) pruned
at the join to two typings---and fans both boxed results into the parent join
through the same replicated input.  Multi-premise cross-products compose at depth
with no starvation.

\subsection{Cost-accounted reactions}
\label{sec:rho-cost}

Reactions are also the natural place to \emph{meter} computation: each
\textsc{comm} is one unit of work, so a resource budget can be charged per
reaction.  \CeTTa{}'s \code{cost} profile does this directly---it reduces an
annotated $\rho$ dialect and reports the resource each reaction consumes.  It is
the executable end of an ongoing program to equip the $\rho$-calculus with a
\emph{cost endofunctor}: a construction that wraps a calculus with resource
accounting once and for all~\cite{costendofunctor}.  We return to it in future
work.

%% ────────────────────────────────────────────────────────────────────
